%%%%%%%%%%%%%%%%%%%%%%%%%%%%%%%%%%%%%%%%%%%%%%%%%%%%%%%%%%%%%%%%%%%%%%%%%%%%%%%%
%  MCM/ICM 2027 -- Paper Skeleton (COMAP-compliant, ready to write in)
%
%  Base       : official COMAP template  corpus/official/MCM-ICM_Summary.tex
%  Fact basis : corpus/official/INDEX.md  (evidence-graded, do not guess rules)
%  Engine     : pdfLaTeX / XeLaTeX / LuaLaTeX.  ALL COMMENTS ARE ASCII ON PURPOSE
%               so the file compiles on every engine; do not paste Chinese into
%               prose -- the report itself must be in English.
%
%  WHAT THIS FILE ALREADY SATISFIES
%    [1] Summary Sheet is PDF page 1, fits on one page, 12pt, English
%    [2] Team control number + page number appear in the header of every page
%    [3] No student / advisor / institution / region name anywhere
%    [4] 25-page cap covers Summary + ToC + body + References + Appendices,
%        so every section below is inside the counted part
%    [5] "Report on Use of AI" is placed AFTER the 25-page part and is excluded
%        from the page count (the only section that is excluded)
%    [6] Section order follows the official content checklist (INDEX.md sec 2.5.3)
%
%  HOW TO USE
%    1. Set \Problem and \Team below.
%    2. Replace every <...> placeholder; delete the placeholder text you replace.
%    3. Compile, then count pages: the part before "Report on Use of AI" must
%       be <= 25 pages.  Trim the appendix first, never the summary.
%    4. Do NOT add a title page, an author block, a running footer, or a
%       "Team Members" line -- each of these breaks the anonymity rule.
%%%%%%%%%%%%%%%%%%%%%%%%%%%%%%%%%%%%%%%%%%%%%%%%%%%%%%%%%%%%%%%%%%%%%%%%%%%%%%%%

\documentclass[12pt]{article}            % >=12pt is a hard COMAP requirement

\usepackage[letterpaper]{geometry}       % US Letter allowed; A4 also allowed
\geometry{left=1in,right=0.75in,top=1in,bottom=1in}

%%%%%%%%%%%%%%%%%%%%%%%%%%%%%%%%%%%%%%%%%%%%%%%%%%%%%%%%%%%%%%%%%%%%%%%%%%%%%%%%
% Team identity -- the ONLY identifying information allowed in the whole file.
% \Problem is the problem letter (MCM: A B C / ICM: D E F).
\newcommand{\Problem}{ABC}               % <-- replace
\newcommand{\Team}{1111111}              % <-- replace with team control number
%%%%%%%%%%%%%%%%%%%%%%%%%%%%%%%%%%%%%%%%%%%%%%%%%%%%%%%%%%%%%%%%%%%%%%%%%%%%%%%%

% ---------------------------- Fonts -----------------------------------------
\usepackage[T1]{fontenc}
\usepackage{newtxtext}                   % Times-like text (official template)
\usepackage{amsmath,amssymb,amsthm}
\usepackage{newtxmath}                   % must come after amsmath

% ---------------------- Layout and typography -------------------------------
\usepackage{graphicx}
\usepackage{xcolor}
\usepackage{booktabs}                    % professional tables
\usepackage{array}
\usepackage{multirow}
\usepackage{siunitx}                     % numbers and units that line up
\usepackage{caption}
\usepackage{subcaption}
\usepackage{float}
\usepackage{enumitem}
\usepackage{listings}                    % code appendix
\usepackage{fancyhdr}
\usepackage[hidelinks]{hyperref}         % load late; keep bookmarks off-page

\sisetup{detect-all,group-separator={,}}

\lstset{
  basicstyle=\ttfamily\footnotesize,
  breaklines=true,
  frame=single,
  numbers=left,
  numberstyle=\tiny,
  captionpos=b,
  showstringspaces=false,
  columns=fullflexible,
  keepspaces=true
}

% ---------------------- Running header (every page) -------------------------
% Hard rule: every page must carry the team control number and the page number
% at the top of the page (instructions.html:1230).  The wording
% "Team # 0000000, Page 6 of 25" given by COMAP is introduced by "for example",
% so the exact string is NOT mandatory -- what is mandatory is that both
% elements are present on every page.
%
% The form below is the official template's own header, unchanged.
% If you prefer the example's exact wording, replace the two lines with:
%     \fancyhead[L]{Team \# \Team, Page \thepage\ of 25}
% and drop the \fancyhead[R] line.
\pagestyle{fancy}
\fancyhf{}
\fancyhead[L]{Team \# \Team}
\fancyhead[R]{Page \thepage}
\renewcommand{\headrulewidth}{0.4pt}
\renewcommand{\footrulewidth}{0pt}

% ---------------------------- Theorems --------------------------------------
\newtheorem{theorem}{Theorem}
\newtheorem{corollary}[theorem]{Corollary}
\newtheorem{lemma}[theorem]{Lemma}
\newtheorem{proposition}[theorem]{Proposition}
\theoremstyle{definition}
\newtheorem{definition}{Definition}
\theoremstyle{remark}
\newtheorem{remark}{Remark}

%%%%%%%%%%%%%%%%%%%%%%%%%%%% PAGE BUDGET (25 pages, hard cap) %%%%%%%%%%%%%%%%%%
% These are planning notes, not content. A workable allocation for 25 pages:
%   Summary Sheet ................ 1   (page 1, must fit on ONE page)
%   Table of Contents ............ 1   (encouraged, COUNTS toward the 25)
%   Restatement + Introduction ... 1.5
%   Assumptions + Justifications . 1.5   <- judges look for justification, not the list
%   Notations .................... 0.5
%   Model Development ............ 7-9   <- the single heaviest section
%   Results ...................... 3-4
%   Testing (sensitivity etc.) ... 2-3
%   Strengths and Weaknesses ..... 1     <- required, do not cut
%   Conclusions .................. 1
%   References ................... 1
%   Appendices ................... 2-4   <- the buffer: cut here first
% Never cut the Testing or Strengths/Weaknesses sections to make room -- they are
% explicit items on the official content checklist.
%%%%%%%%%%%%%%%%%%%%%%%%%%%%%%%%%%%%%%%%%%%%%%%%%%%%%%%%%%%%%%%%%%%%%%%%%%%%%%%

\begin{document}

% ==============================================================================
%  PAGE 1 -- SUMMARY SHEET.  Must be the first page of the PDF and must not
%  spill onto a second page.  Write this LAST: COMAP states judges weight the
%  summary heavily and that a weak summary is usually the end of the paper's
%  chances. Do not restate the problem at length and do not paste boilerplate
%  from your Introduction.
% ==============================================================================
\graphicspath{{.}}                       % keep figures beside this .tex file
\DeclareGraphicsExtensions{.pdf,.jpg,.tif,.png}
\thispagestyle{fancy}                    % keep the header, so page 1 is numbered too
\vspace*{-16ex}

\centerline{\begin{tabular}{*3{c}}
	\parbox[t]{0.3\linewidth}{\begin{center}\textbf{Problem Chosen}\\ \Large \textcolor{red}{\Problem}\end{center}}
	& \parbox[t]{0.3\linewidth}{\begin{center}\textbf{2027\\ MCM/ICM\\ Summary Sheet}\end{center}}
	& \parbox[t]{0.3\linewidth}{\begin{center}\textbf{Team Control Number}\\ \Large \textcolor{red}{\Team}\end{center}}	\\
	\hline
\end{tabular}}

%%%%%%%%%%% Begin Summary %%%%%%%%%%%
% Delete the placeholders below and write the summary here. Target shape:
%   - a bold lead sentence stating what you did and the headline number
%   - one paragraph per model / stage, naming the method and saying WHY
%   - concrete results with numbers and units, never adjectives alone
%   - one sentence on how you tested the model
%   - one sentence of conclusion / recommendation
% No figures, no tables, no citations in the summary.

\begin{center}
\textbf{<One sentence: what we were asked, what we built, and the headline result>}
\end{center}

<Paragraph 1 -- the problem in your own words and why it matters, closing with a
one-sentence statement of your overall approach.>

<Paragraph 2 -- the first model: name it, say what it represents, and give the
justification for choosing it over the obvious alternative.>

<Paragraph 3 -- the second model / solution stage, including how the two stages
connect.>

<Paragraph 4 -- the results, with numbers, units, and the comparison or baseline
they should be read against.>

<Paragraph 5 -- how you validated the model: sensitivity, robustness, error
analysis; state the range over which the conclusion holds.>

<Paragraph 6 -- the conclusion and its practical implication for the
decision-maker in the problem statement.>

\vspace{1ex}
\noindent\textbf{Keywords:} <keyword>; <keyword>; <keyword>; <keyword>
%%%%%%%%%%% End Summary %%%%%%%%%%%

%%%%%%%%%%%%%%%%%%%%%%%%%%%%%%
\clearpage
\pagestyle{fancy}
% ==============================================================================
%  PAGE 2 (counted) -- TABLE OF CONTENTS.  Encouraged by COMAP and it counts
%  toward the 25 pages.  Keep it to one page; it is generated automatically.
% ==============================================================================
\tableofcontents
\newpage
% NOTE: the official template resets the page counter here so that the body
% starts at "Page 1" while the summary sheet is unnumbered.  We deliberately do
% NOT reset it: the 25-page cap covers the summary sheet too, so numbering the
% submission continuously (summary = page 1) keeps "Page N" honest against the
% cap.  If your team prefers the official flow, add \setcounter{page}{1} here.

% ==============================================================================
%  BODY -- 25-page counted region starts here.
%  Official content checklist (each item MUST have a home in what follows):
%    ToC / restatement / clear statement of variables and assumptions /
%    justification of the assumptions / analysis that motivates the model /
%    derivation summary in the body with long algebra pushed to the appendix /
%    model design AND how it was tested (error, sensitivity, stability) /
%    strengths and weaknesses / conclusions with results stated explicitly /
%    references.  All ten appear below -- keep them all.
% ==============================================================================

\section{Introduction}
% Answers: what is the problem, why does it matter, what has been done before,
% and what does THIS paper do that is different?
\subsection{Background and Restatement of the Problem}
% Restate the problem in your own words, in your own structure. Do not copy the
% problem statement. End with the specific questions you will answer, numbered
% or bulleted, so the reader can map results back to them.
<...>

\subsection{Literature and Data Review}
% Cite real, checked sources with inline citations. Every borrowed figure, data
% table, idea, or image needs a citation here AND in the References section;
% uncited borrowed material is the most common route to a plagiarism finding.
<...>

\subsection{Our Approach}
% One paragraph, plus a flow diagram if you have one. Say what the stages are
% and how they feed each other.
<...>

\section{Assumptions and Justifications}
% Official checklist item. A bare list of assumptions earns nothing: each one
% needs a stated justification, and ideally a note on what happens if it fails.
% Keep the ones you actually use in the model; drop the decorative ones.
\begin{enumerate}[label=\textbf{A\arabic*.},leftmargin=*]
  \item \textbf{<Assumption.>} <Justification -- cite data or a source where you can.>
  \item \textbf{<Assumption.>} <Justification.>
  \item \textbf{<Assumption.>} <Justification.>
\end{enumerate}

\section{Notations}
% Official checklist item: a clear statement of variables. Use a table; define
% every symbol that appears in a later equation. Keep units in a column.
\begin{table}[H]
\centering
\caption{Symbols used in this paper.}
\label{tab:notations}
\begin{tabular}{cll}
\toprule
Symbol & Definition & Unit \\
\midrule
$<x>$  & $<$...$>$ & $<$...$>$ \\
$<y>$  & $<$...$>$ & $<$...$>$ \\
\bottomrule
\end{tabular}
\end{table}

\section{Model Development}
% The heaviest section. The official checklist asks for the ANALYSIS that
% motivates the model before the model itself -- show why this model and not
% another. Keep long algebra out of the body: state the derivation and push the
% steps to Appendix A. Derivation/calculation summary belongs HERE, detail in
% the appendix.
\subsection{Model Selection and Motivation}
% State each candidate approach you considered and the criterion that decided
% between them (data availability, tractability, interpretability, ...).
<...>

\subsection{<Name of your first model>}
\begin{equation}
  \label{eq:main}
  <...>
\end{equation}
% Explain each term of Eq. (\ref{eq:main}) in words, then state how it is solved
% (closed form, numerical scheme, solver and version).
<...>

\subsection{<Second model / stage, if any>}
<...>

\section{Results}
% Report numbers, not impressions. Every figure and table needs a caption that
% says what it shows and what the reader should conclude from it; refer to each
% one from the text.
\begin{figure}[H]
  \centering
  % \includegraphics[width=0.8\linewidth]{fig1}
  \caption{<What is plotted; what the reader should take from it.>}
  \label{fig:result1}
\end{figure}

\section{Model Testing}
% Official checklist item, and one of the categories in COMAP's published sample
% rubric, so it carries real weight. All three sub-parts need something in them.
\subsection{Sensitivity Analysis}
% Vary the parameters that matter and report the range over which your
% conclusion survives. Give the range, not just "the result is robust".
<...>

\subsection{Stability and Robustness}
% Re-run under a different discretisation, solver setting, or resampled data.
<...>

\subsection{Error Analysis}
% Sources of error (measurement, model structure, numerical), their size, and
% which direction they push the answer.
<...>

\section{Strengths and Weaknesses}
% Official checklist item: discuss the obvious strengths AND weaknesses of your
% model or approach. A weakness named with an honest bound reads as competence;
% claiming there are none does not.
\subsection{Strengths}
<...>
\subsection{Weaknesses and Limitations}
<...>

\section{Conclusions}
% Official checklist item: state the conclusions AND report the results
% explicitly. Answer, in order, each question you posed in Section 1.
<...>

% ==============================================================================
%  REFERENCES -- counted toward the 25 pages.
%  Cite in the text as well as here. Note that AI tools you used must ALSO be
%  listed here (see the Report on Use of AI at the end of this file).
% ==============================================================================
\begin{thebibliography}{9}
\bibitem{<key1>} <Author>. \textit{<Title>}. <Publisher / Journal>, <Year>.
\bibitem{<key2>} <Author>. <Title>. <URL>, accessed <date>.
\end{thebibliography}

\appendix
% ==============================================================================
%  APPENDICES -- counted toward the 25 pages. This is the buffer: when you are
%  over the limit, cut here before cutting Testing or Strengths/Weaknesses.
%  Do not put the only copy of a derivation here; summarise it in the body too.
% ==============================================================================
\section{Derivations Omitted from the Body}
<...>

\section{Additional Tables and Figures}
<...>

\section{Code}
% COMAP will not run your code, and non-paper files are not accepted -- but the
% code listing here is what makes the work reproducible. Keep it short; long
% listings eat the page budget.
% \lstinputlisting[language=Python,caption={<model.py>}]{model.py}
<...>

% ==============================================================================
%  REPORT ON USE OF AI -- OUTSIDE the 25-page solution.
%  Required whenever ANY AI-assisted tool was used, which includes LLMs,
%  generative AI, MATHEMATICS SOFTWARE, TRANSLATION SOFTWARE, and code
%  completion. It goes after the 25 pages, in the SAME PDF file, has no page
%  limit, and does NOT count toward the 25.
%
%  Two disclosures are required, not one: (a) inline citations in the body plus
%  entries in the References section above, and (b) this section. Teams that
%  skip (a) are the ones at risk -- undisclosed AI use is what COMAP acts on.
%
%  State, for each tool: name, version/date, model name where the tool has one,
%  what it was used for, and how you verified the output. For translation-only
%  use a single declarative sentence is enough; for tools with no discrete
%  query (translation software, inline code completion) a sentence is also
%  enough. For chat-style tools give the exact query and the complete output.
% ==============================================================================
\clearpage
\fancyhead[R]{Page \thepage}             % outside the capped count: drop "of 25"
\section*{Report on Use of AI}
\addcontentsline{toc}{section}{Report on Use of AI (not counted toward the 25-page limit)}

\begin{enumerate}[leftmargin=*]
  \item <Tool name> (<version/date>, <model name if the tool has one>)
        \begin{itemize}[nosep]
          \item Purpose: <what it was used for -- e.g. language polishing of
                Sections 2 and 5, or translation of the team's draft into English.>
          \item Query: <the exact wording entered into the tool.>
          \item Output: <the complete output returned.>
          \item Verification: <how the team checked it -- compared against
                <source>, recomputed independently, corrected <what>.>
        \end{itemize}
  \item <Tool with no discrete query, e.g. a translation tool or a code
        completion tool> (<version/date>): <one declarative sentence stating
        what it was used for.>
\end{enumerate}

\noindent <If no AI tool of any kind was used, say so explicitly rather than
leaving this section out.>

%%%%%%%%%%%%%%%%%%%%%%%%%%%%%%
\end{document}
